% answers_short.tex —— 2020 年·全国Ⅱ卷·理科 · 简短答案速查
% 由 tools/make_short_answers.py 自 answers.tex 的 \gksolhead/\gkanswer 提取，
% 并与 2020/2020_数学真题_全国2理_答案要点.md 逐题比对后生成，勿手改（改源 MD 后重跑）。

\begin{gknotice}
\textbf{使用说明}：本节只给最终答案便于快速核对，完整推导见解析版。答案为编写者独立推导，\textbf{非官方标准答案}。
\end{gknotice}

\gkpart{选择题}
\begin{center}\begin{tabular}{ccccccccccccc}
\toprule
题号 & 1 & 2 & 3 & 4 & 5 & 6 & 7 & 8 & 9 & 10 & 11 & 12 \\
\midrule
答案 & A & D & B & C & B & C & A & B & D & C & A & C \\
\bottomrule
\end{tabular}\end{center}

\gkpart{填空题}
\begin{itemize}[leftmargin=2.6em]
  \item[\textbf{13.}] $\dfrac{\sqrt2}2$。由 $(k\boldsymbol a-\boldsymbol b)\cdot\boldsymbol a=0$ 得 $k=\boldsymbol a\cdot\boldsymbol b=\cos45^\circ$。
  \item[\textbf{14.}] $36$。先分组后分配：$\dfrac{\mathrm C_4^2\mathrm C_2^1\mathrm C_1^1}{2!}\times3!=6\times6=36$。
  \item[\textbf{15.}] $2\sqrt3$。由 $|z_1+z_2|^2+|z_1-z_2|^2=2(|z_1|^2+|z_2|^2)=16$ 得 $|z_1-z_2|^2=16-4=12$。
  \item[\textbf{16.}] ①③④。$p_1$、$p_4$ 为真，$p_2$、$p_3$ 为假，代入四种复合形式即得。
\end{itemize}

\gkpart{解答题}
\begin{itemize}[leftmargin=2.6em]
  \item[\textbf{17.}] (1) $A=\dfrac{2\pi}3$。由正弦定理化边得 $a^2-b^2-c^2=bc$，再用余弦定理得 $\cos A=-\dfrac12$；(2) $3+2\sqrt3$。由 $a=3$、$A=\dfrac{2\pi}3$ 得 $b^2+c^2+bc=9$，而 $(b+c)^2=9+bc\leqslant9+\dfrac{(b+c)^2}4$，故 $b+c\leqslant2\sqrt3$
  \item[\textbf{18.}] (1) $12000$。样本平均数 $\bar y=\dfrac{1200}{20}=60$，乘地块数 $200$ 得 $60\times200=12000$；(2) $0.94$。$r=\dfrac{800}{\sqrt{80\times9000}}=\dfrac{2\sqrt2}3\approx0.94$；(3) 分层抽样（按植物覆盖面积分层）。地块间差异大而层内差异小，分层后样本结构与总体一致，估计更准确
  \item[\textbf{19.}] (1) $e=\dfrac12$。由 $|CD|=\dfrac43|AB|$ 得 $4c=\dfrac43\cdot\dfrac{2b^2}a$，即 $3ac=2(a^2-c^2)$，解 $2e^2+3e-2=0$；(2) $C_1:\dfrac{x^2}{36}+\dfrac{y^2}{27}=1$，$C_2:y^2=12x$。由 $|MF|=x_M+c=5$ 与两曲线方程联立得 $c=3$，故 $a=6$，$b^2=27$
  \item[\textbf{20.}] (1) 证明。$MN\parallel BB_1\parallel AA_1$；又 $BC\perp AM$、$BC\perp MN$，故 $BC\perp$ 平面 $A_1AMN$，而 $B_1C_1\parallel BC$ 且 $B_1C_1\subset$ 平面 $EB_1C_1F$，故两平面垂直；(2) $\dfrac{\sqrt{10}}{10}$。以 $M$ 为原点、$MA$ 为 $y$ 轴、$MB$ 为 $x$ 轴建系，由 $AO\parallel$ 平面 $EB_1C_1F$ 与 $AO=AB$ 得 $|B_1E|=\dfrac{\sqrt{10}}3AB$、$B_1E$ 的横向分量为 $\dfrac13AB$
  \item[\textbf{21.}] (1) 在 $\left(0,\dfrac\pi3\right)$、$\left(\dfrac{2\pi}3,\pi\right)$ 单调递增，在 $\left(\dfrac\pi3,\dfrac{2\pi}3\right)$ 单调递减。由 $f'(x)=2\sin^2x(4\cos^2x-1)$ 定号；(2) 证明。$f$ 在 $[0,\pi]$ 上的最大值为 $f\left(\dfrac\pi3\right)=\dfrac{3\sqrt3}8$，最小值为 $f\left(\dfrac{2\pi}3\right)=-\dfrac{3\sqrt3}8$，且 $f$ 以 $\pi$ 为周期；(3) 证明。记 $b_k=|\sin2^kx|$，由（2）得 $b_k^2b_{k+1}\leqslant\dfrac{3\sqrt3}8$（$k=0,\cdots,n-1$），连乘后与 $\prod b_k^3$ 比较，得 $\prod b_k^2\leqslant\left(\dfrac{3\sqrt3}8\right)^{2n/3}=\left(\dfrac34\right)^n$
\end{itemize}

\gkpart{选做题}
\begin{itemize}[leftmargin=2.6em]
  \item[\textbf{22.}] (1) $C_1:x+y=4$（$0\leqslant x\leqslant4$），$C_2:x^2-y^2=4$。前者由 $4\cos^2\theta+4\sin^2\theta=4$ 消参，后者由 $x^2-y^2=4$ 消参；(2) $\rho=\dfrac{17}5\cos\theta$。交点 $P\left(\dfrac52,\dfrac32\right)$，圆心 $\left(\dfrac{17}{10},0\right)$，故 $\rho=2\cdot\dfrac{17}{10}\cos\theta$
  \item[\textbf{23.}] (1) $\left(-\infty,\dfrac32\right]\cup\left[\dfrac{11}2,+\infty\right)$。$a=2$ 时 $f(x)=|x-4|+|x-3|$，分段求解 $f(x)\geqslant4$；(2) $a\leqslant-1$ 或 $a\geqslant3$。$f(x)$ 的最小值为 $(a-1)^2$，故 $(a-1)^2\geqslant4$
\end{itemize}
